% Gabarit pandoc des corrigés publiés (brut et guidé). Contenu sans marque ni nom de site.
% Compilation : lualatex (latexmk -lualatex).
% Le Markdown reste la source unique : ne pas éditer ce fichier .tex à la main.
\documentclass[11pt,a4paper]{article}

% --- Moteur, polices (Latin Modern) et langue ---------------------------------
\usepackage{iftex}
\ifLuaTeX\else\errmessage{Ce document se compile avec lualatex}\fi
\usepackage{amsmath,amssymb,mathtools}
\usepackage{fontspec}
% mathrm=sym : \mathrm prend ses glyphes dans la police mathématique (ℓ, μ, Ω), comme KaTeX.
\usepackage[math-style=TeX,bold-style=TeX,mathrm=sym]{unicode-math}
% Repli pour les glyphes absents de Latin Modern (indices et exposants Unicode,
% flèches, filets des schémas en texte, symboles) : DejaVu, puis Latin Modern Math
% (ℓ, lettres mathématiques Unicode des légendes).
\directlua{
  luaotfload.add_fallback("repliTexte", {"DejaVuSerif:mode=harf;", "DejaVuSans:mode=harf;", "file:latinmodern-math.otf:mode=harf;"})
  luaotfload.add_fallback("repliMono", {"DejaVuSansMono:mode=harf;", "DejaVuSans:mode=harf;", "file:latinmodern-math.otf:mode=harf;"})
}
\setmainfont{Latin Modern Roman}[RawFeature={fallback=repliTexte}]
\setsansfont{Latin Modern Sans}[RawFeature={fallback=repliTexte}]
\setmonofont{Latin Modern Mono}[RawFeature={fallback=repliMono}]
\setmathfont{Latin Modern Math}
\usepackage[french]{babel}
\frenchsetup{SmallCapsFigTabCaptions=false}

% --- Mise en page --------------------------------------------------------------
\usepackage[a4paper,margin=2.2cm,headheight=14pt,footskip=1.2cm]{geometry}
\usepackage{microtype}
\setlength{\emergencystretch}{4em}
\tolerance=1500
\setlength{\parindent}{0pt}
\setlength{\parskip}{0.55\baselineskip plus 2pt minus 1pt}
\usepackage{xcolor}
% Gris nommés : babel-french rend « ! » actif dans le document.
\definecolor{gris2}{gray}{0.98}
\definecolor{gris3}{gray}{0.97}
\definecolor{gris5}{gray}{0.95}
\definecolor{gris35}{gray}{0.65}
\definecolor{gris40}{gray}{0.60}
\definecolor{gris45}{gray}{0.55}
\definecolor{gris50}{gray}{0.50}
\definecolor{gris55}{gray}{0.45}
\definecolor{gris60}{gray}{0.40}
\definecolor{gris65}{gray}{0.35}
\definecolor{gris70}{gray}{0.30}
\definecolor{gris80}{gray}{0.20}
\definecolor{gris85}{gray}{0.15}
\usepackage{graphicx}
\usepackage{float}
\usepackage{caption}
\captionsetup{font=small,labelformat=empty,justification=centering}
\usepackage{booktabs,longtable,array,multirow,calc}
\usepackage{siunitx}
\sisetup{locale=FR}
\usepackage[version=4]{mhchem}
\usepackage{listings}
\usepackage{fancyhdr}
\usepackage{lastpage}
\usepackage[most]{tcolorbox}
\usepackage{enumitem}
\setlist{nosep,leftmargin=*,topsep=2pt}
\usepackage[hidelinks,unicode]{hyperref}
\usepackage{bookmark}
\hypersetup{pdftitle={Mines-Ponts MP-MPI-PC-PSI — Français 2026},pdfauthor={Rédigé par IA},pdfcreator={LaTeX},pdfsubject={Corrigé brut}}
\usepackage{xurl}

% --- Pandoc : listes, tableaux, figures ----------------------------------------
\providecommand{\passthrough}[1]{#1}
\providecommand{\tightlist}{\setlength{\itemsep}{0pt}\setlength{\parskip}{0pt}}

\makeatletter
\def\maxwidth{\ifdim\Gin@nat@width>\linewidth\linewidth\else\Gin@nat@width\fi}
\def\maxheight{\ifdim\Gin@nat@height>0.55\textheight 0.55\textheight\else\Gin@nat@height\fi}
\makeatother
\setkeys{Gin}{width=\maxwidth,height=\maxheight,keepaspectratio}
\setcounter{secnumdepth}{-\maxdimen}
\providecommand{\pandocbounded}[1]{#1}

% --- Code (listings) -----------------------------------------------------------
\lstset{
  basicstyle=\ttfamily\small,
  columns=fullflexible,
  keepspaces=true,
  breaklines=true,
  breakatwhitespace=false,
  postbreak=\mbox{\textcolor{gris50}{$\hookrightarrow$}\space},
  frame=single,
  rulecolor=\color{gris35},
  framesep=4pt,
  xleftmargin=5pt,xrightmargin=5pt,
  backgroundcolor=\color{gris3},
  keywordstyle=\bfseries,
  commentstyle=\color{gris65},
  stringstyle=\color{gris80},
  showstringspaces=false,
  tabsize=4,
  extendedchars=true,
  aboveskip=6pt,belowskip=6pt,
}

% --- Macros des corrigés (KaTeX ↔ LaTeX) ---------------------------------------
% Les corrigés sont écrits pour KaTeX : on définit ici ce que KaTeX connaît et
% que LaTeX n'a pas, sans jamais retoucher le Markdown.
\providecommand{\R}{\mathbb{R}}
\providecommand{\N}{\mathbb{N}}
\providecommand{\Z}{\mathbb{Z}}
\providecommand{\Q}{\mathbb{Q}}
\providecommand{\C}{\mathbb{C}}
\providecommand{\dd}{\mathrm{d}}
\providecommand{\llbracket}{\lBrack}
\providecommand{\rrbracket}{\rBrack}
\providecommand{\lBrace}{\lbrace}
\providecommand{\rBrace}{\rbrace}
\providecommand{\Box}{\square}
\providecommand{\Diamond}{\lozenge}
\providecommand{\varnothing}{\emptyset}
\providecommand{\argmax}{\operatorname*{arg\,max}}
\providecommand{\argmin}{\operatorname*{arg\,min}}
\renewcommand{\leq}{\leqslant}
\renewcommand{\geq}{\geqslant}

% Équation hors texte : réduite à la largeur de la ligne si elle déborde.
\newsavebox{\boiteequation}
\newcommand{\equationajustee}[1]{%
  \sbox{\boiteequation}{$\displaystyle #1$}%
  \ifdim\wd\boiteequation>\linewidth
    \[\resizebox{\linewidth}{!}{\usebox{\boiteequation}}\]%
  \else
    \[#1\]%
  \fi}

% --- Encadrés pédagogiques (sobres, lisibles en noir et blanc) -----------------
\tcbset{
  encadre/.style={
    enhanced,breakable,sharp corners,
    colback=white,colframe=gris55,
    boxrule=0.4pt,leftrule=2.2pt,
    left=6pt,right=6pt,top=3pt,bottom=3pt,
    before skip=6pt,after skip=6pt,
    fonttitle=\bfseries\small,coltitle=black,
    attach title to upper,after title={\quad},
    parbox=false,
  }
}
\newtcolorbox{indication}[1]{encadre,colframe=gris45,colback=gris2,title={Indication~#1}}
\newtcolorbox{methode}{encadre,colframe=gris70,title={Méthode}}
\newtcolorbox{piege}{encadre,colframe=gris85,boxrule=0.8pt,leftrule=2.2pt,title={Piège}}
\newtcolorbox{jury}{encadre,colframe=gris60,colback=gris5,title={Jury}}
\newtcolorbox{encadrenote}{encadre,colframe=gris40}

% --- Titres de section ---------------------------------------------------------
\usepackage{titlesec}
\titleformat{\section}{\Large\bfseries}{}{0pt}{}
\titleformat{\subsection}{\large\bfseries}{}{0pt}{}
\titleformat{\subsubsection}{\normalsize\bfseries}{}{0pt}{}
\titlespacing*{\section}{0pt}{14pt plus 4pt minus 2pt}{6pt}
\titlespacing*{\subsection}{0pt}{10pt plus 3pt minus 2pt}{4pt}
\titlespacing*{\subsubsection}{0pt}{8pt plus 2pt minus 2pt}{2pt}

% --- En-tête et pied de page ---------------------------------------------------
\pagestyle{fancy}
\fancyhf{}
\fancyhead[L]{\small\textsc{Corrigé brut} · Mines-Ponts MP-MPI-PC-PSI}
\fancyhead[R]{\small Français 2026}
\fancyfoot[L]{\footnotesize Corrigé rédigé par IA, non officiel}
\fancyfoot[R]{\footnotesize \thepage\,/\,\pageref*{LastPage}}
\renewcommand{\headrulewidth}{0.4pt}
\renewcommand{\footrulewidth}{0pt}
\fancypagestyle{plain}{\fancyhf{}\fancyfoot[R]{\footnotesize \thepage\,/\,\pageref*{LastPage}}\fancyfoot[L]{\footnotesize Corrigé rédigé par IA, non officiel}\renewcommand{\headrulewidth}{0pt}}

\begin{document}
\thispagestyle{plain}
\begin{center}
  {\small\scshape Corrigé brut}\par\smallskip
  {\LARGE\bfseries Mines-Ponts MP-MPI-PC-PSI — Français 2026\par}\medskip
  {\large Mines-Ponts · MP-MPI-PC-PSI · Français · 2026\par}
\end{center}
\vspace{2pt}\hrule\vspace{8pt}

\emph{Corrigé rédigé par IA, non officiel. Signalez toute erreur.}

\textbf{Problématique} : l'expérience de la nature, telle que la donnent à lire les trois œuvres, est-elle cette conquête dont l'homme rapporterait un mot qui le rendrait maître de son être, ou bien ne vérifie-t-elle la phrase de Sand qu'à condition d'en renverser le vocabulaire de la prise ?

\begin{center}\rule{0.5\linewidth}{0.5pt}\end{center}

En 1872, dans un article du \emph{Temps} consacré à la forêt de Fontainebleau, George Sand ne décrit pas un paysage : elle plaide pour une forêt, à une époque où ses futaies sont exploitées et où les peintres ont obtenu que certaines parcelles soient soustraites à la coupe. Le ton est celui d'une prophétesse plus que d'une promeneuse : « Quelle conquête à entreprendre pour l'homme, et je dis pour tout homme actuellement vivant ou à naître ! Entrer dans la nature, chercher l'oracle de la forêt sacrée et rapporter le mot, ne fût-ce qu'un mot, qui doit répandre sur toute sa vie le charme profond de la possession de son être ! » La phrase dessine un trajet en trois gestes, entrer, chercher, rapporter, et le couronne par un état, la possession de soi. Deux lexiques s'y croisent : celui de la conquête, qui est celui de la prise, et celui de l'oracle, qui est celui du sacré. Mais le plus remarquable est le déplacement qu'elle opère en chemin : la conquête annoncée n'a pas pour objet la nature ; ce que l'homme doit finalement posséder, c'est « son être ». Encore ce butin tient-il dans « un mot », et la restrictive « ne fût-ce qu'un mot » tempère d'avance l'enthousiasme de l'exclamation.

Ce programme suppose pourtant beaucoup : qu'il existe un seuil de la nature, que l'homme se tienne d'ordinaire au dehors et puisse y entrer ; que la nature ait quelque chose à dire ; que celui qui l'a entendue puisse revenir et transmettre. Or les trois œuvres au programme racontent ou pensent des immersions. Dans \emph{Vingt mille lieues sous les mers}, Jules Verne fait descendre le professeur Aronnax dans les profondeurs à bord du \emph{Nautilus} du capitaine Nemo ; dans \emph{Le Mur invisible}, Marlen Haushofer enferme sa narratrice, seule avec quelques bêtes, dans une vallée des Alpes autrichiennes que ferme une paroi transparente ; dans \emph{La Connaissance de la vie}, Georges Canguilhem se demande à quelles conditions la pensée peut connaître le vivant sans le manquer. On se demandera donc dans quelle mesure l'expérience de la nature, telle que ces œuvres la donnent à lire, est la conquête dont l'homme reviendrait en possession de lui-même, ou si elle ne vérifie la phrase de Sand qu'à condition d'en renverser le vocabulaire.

Nous verrons d'abord que les trois œuvres suivent bien le trajet que trace Sand : on ne connaît la nature qu'en y entrant, en l'interrogeant et en écrivant ce qu'on en a appris. Nous montrerons ensuite qu'elles récusent les mots par lesquels Sand décrit ce trajet : la conquête y tourne à la prédation, l'oracle y reste muet, le retour y est incertain. Nous établirons enfin que la formule retrouve sa justesse si on la lit à partir de son dernier terme : la possession de son être n'est pas une prise, mais la capacité, pour un vivant parmi d'autres, de se donner des normes et de répondre de ce qui l'entoure.

\begin{center}\rule{0.5\linewidth}{0.5pt}\end{center}

Le premier geste que réclame Sand, entrer, est celui dont les trois œuvres font la condition de toute connaissance. Aronnax est d'abord un savant de cabinet : professeur suppléant au Muséum de Paris, il a publié un ouvrage en deux volumes sur les mystères des grands fonds sous-marins, qu'il n'a jamais vus. Le roman consiste à faire passer ce savoir du livre à la chose, et Verne donne à cette entrée une figure littérale : la promenade en scaphandre dans la forêt sous-marine de l'île Crespo. C'est bien une forêt, et le naturaliste y entre à pied ; il y voit de ses yeux ce que ses livres ne lui avaient pas appris, une végétation qui monte droit vers la surface. Canguilhem formule la même exigence dans l'ordre de la méthode. L'étude consacrée à l'expérimentation en biologie animale montre que le vivant ne se laisse pas traiter comme un corps physique : il est un individu, une totalité, une histoire, et l'on ne peut isoler une fonction, ni conclure d'une espèce à une autre, sans altérer ce qu'on prétend connaître. Le laboratoire lui-même est un milieu, avec ses contraintes propres. Pour connaître le vivant, il faut donc entrer dans les conditions de sa vie au lieu de l'en extraire. Chez Haushofer, enfin, l'entrée n'est pas choisie, elle est subie : venue passer quelques jours au pavillon de chasse de sa cousine Luise et de son mari Hugo, la narratrice découvre au matin que ses hôtes ne sont pas revenus du village et qu'une paroi invisible ferme la route. Le seuil de Sand est devenu un mur qui ne se franchit que dans un sens, mais l'entrée est totale : une citadine doit apprendre la nature par le corps.

Cette entrée ouvre une recherche, et c'est le second geste de Sand : chercher l'oracle, c'est-à-dire interroger une nature qui ne livre pas d'elle-même son sens. Le roman de Verne s'ouvre sur une énigme, ce prétendu monstre marin que les navires croisent en 1866, et il en multiplie d'autres, des ruines de l'Atlantide aux poulpes géants que la science tenait pour fabuleux. Il oppose surtout deux manières de connaître : Conseil, le domestique flamand, range tout poisson dans sa classe et son ordre sans savoir le reconnaître, tandis que Ned Land, le harponneur, le reconnaît sans savoir le classer ; Aronnax note qu'à eux deux ils feraient un naturaliste. Interroger la nature, c'est ajuster la nomenclature à l'expérience, et d'abord chercher le mot juste : quand Ned demande ce qu'est une perle, Aronnax lui donne tour à tour le mot du poète, celui du chimiste et celui du naturaliste, comme si aucun nom ne suffisait seul à dire la chose. Canguilhem donne à cette recherche sa portée philosophique : le vivant pose à la pensée un problème que ne pose pas la matière inerte, parce qu'il est lui-même un centre d'activité qui organise son milieu. Le chercheur ne trouve donc pas seulement des faits, il rencontre des allures, des préférences, des valeurs. La narratrice d'Haushofer, elle, apprend à déchiffrer des signes vitaux : l'état du ciel avant la fenaison, l'humeur de sa vache Bella, le comportement de son chien Lynx, l'approche de la neige. Son oracle est quotidien, mais il est bien une parole à interpréter, et sa survie dépend de la justesse de l'interprétation.

Reste à « rapporter le mot », et les trois œuvres sont, à la lettre, des rapports. Le roman de Verne se donne pour le récit qu'Aronnax rédige après son évasion ; \emph{Le Mur invisible} est le texte que la narratrice écrit sur ce qu'elle appelle elle-même son « rapport », pour ne pas perdre la raison ; \emph{La Connaissance de la vie} est ce qu'un philosophe rapporte de sa fréquentation des biologistes. Et ce mot rapporté revient bien, comme le veut Sand, sur celui qui le rapporte. La narratrice, dont le texte ne donne jamais le nom, se découvre étrangère à la femme qu'elle était à la ville, à ses soucis et à son apparence ; Aronnax, embarqué pour détruire un monstre, finit partagé entre l'envie de fuir et le désir de poursuivre l'exploration. Canguilhem fait de ce retour sur soi la condition même du savoir biologique : la pensée qui connaît la vie est le fait d'un vivant, et, selon la formule de son introduction, « la pensée du vivant doit tenir du vivant l'idée du vivant ». Connaître la vie, c'est aussi se connaître comme vivant.

\begin{center}\rule{0.5\linewidth}{0.5pt}\end{center}

Les trois œuvres suivent donc le trajet de Sand. Mais le suivent-elles dans les termes où elle le décrit ? Aronnax est un prisonnier, la narratrice une captive, l'expérimentateur un homme qui risque d'abîmer ce qu'il veut connaître : c'est le vocabulaire de la phrase, plus que son mouvement, que les œuvres mettent en cause.

Sand parle de « conquête » et de « possession ». Elle ne prône pas la domination de la nature, puisque sa conquête a pour objet « son être » ; mais le mot vient de la guerre, et l'on entend sous lui le programme cartésien qui invitait les hommes à se rendre « comme maîtres et possesseurs de la nature ». Or ce vocabulaire de la prise, dès qu'il s'applique à la nature, les œuvres l'instruisent à charge. Le geste conquérant par excellence est chez Verne aussitôt démenti : lorsque Nemo plante au pôle Sud son pavillon noir et prend possession, en son propre nom, de terres qu'aucun homme n'habite, le \emph{Nautilus} se trouve peu après bloqué sous la banquise par le basculement d'un iceberg, et l'équipage manque mourir asphyxié. Le roman est aussi un livre de tueries : Nemo, qui interdit à Ned Land de harponner des baleines inoffensives pour le seul plaisir et lui reproche les ravages des baleiniers, massacre peu après une troupe de cachalots qu'il juge nuisibles ; et le proscrit qui défendait les opprimés finit par éperonner un navire de guerre et regarder couler son équipage. La conquête rend Nemo semblable à ce qu'il hait. Chez Haushofer, le procès est plus radical encore. Le seul être humain vivant que rencontre la narratrice est un homme qui surgit à l'alpage et tue à coups de hache le jeune taureau et le chien ; elle l'abat d'un coup de fusil. Cet inconnu sans nom ni parole est la prise à l'état pur, détruire sans raison. Le roman ne propose pas pour autant une innocence : la narratrice tue elle aussi, chasse le gibier pour manger et le fait à contrecœur. Canguilhem, enfin, ôte à la maîtrise son fondement. Dans « Machine et organisme », il renverse le modèle cartésien : loin que l'organisme s'explique par la machine, c'est la machine qui s'explique par l'organisme, la technique prolongeant une activité vitale. On ne possède pas comme un objet démontable ce dont on procède.

Sand suppose ensuite une « forêt sacrée » qui rend des oracles. Or la nature, dans les trois œuvres, ne parle pas. \emph{Le Mur invisible} est un roman de montagne qui refuse le lyrisme de la montagne : la beauté d'un matin d'été ne console de rien, des bêtes meurent sans raison, la corneille blanche que la narratrice observe est rejetée par les autres corneilles parce qu'elle est blanche. Et le mur, événement central du livre, n'est jamais expliqué : ni cause, ni auteur, ni sens. Canguilhem donne à ce silence son explication théorique. Dans « Le vivant et son milieu », il reprend à von Uexküll la distinction entre le milieu propre d'un animal, fait des seuls signaux qui ont pour lui une signification, et l'environnement que décrit le physicien : la tique ne vit pas dans le même monde que le savant qui l'observe. Il n'y a donc pas de sens en soi de la nature, qu'un oracle délivrerait ; le sens naît du rapport qu'un vivant institue avec ce qui l'entoure. Chez Verne, enfin, le seul « mot » que le voyage fasse apparaître sous la mer est humain : dans les ruines de l'Atlantide, Nemo trace sur un rocher, avec un morceau de pierre crayeuse, le mot « ATLANTIDE ». Ce n'est pas la mer qui parle, c'est un homme qui écrit sur elle. Et le livre se clôt sur la question de l'Ecclésiaste, demandant qui a jamais pu sonder les profondeurs de l'abîme : l'oracle n'a pas répondu.

Le troisième geste, « rapporter », est peut-être le plus fragile, car il suppose qu'on revienne et qu'on ait quelqu'un à qui parler. La narratrice d'Haushofer ne revient pas : elle écrit au pavillon un rapport que personne ne lira peut-être, et qui s'arrête non parce que son histoire est achevée, mais parce qu'elle n'a plus de papier. Aronnax revient, mais il sait que son récit est invérifiable et se demande lui-même si on le croira. Quant à Nemo, qui détient le savoir le plus complet du voyage, il ne rapporte rien aux hommes : il a confié l'histoire de sa vie à un manuscrit que le dernier survivant de l'équipage doit jeter à la mer, et le roman s'achève sans que l'on sache si le \emph{Nautilus} a résisté au Maelström. Le mot reste au fond. Canguilhem, enfin, met en doute le « charme » même du mot rapporté : son introduction rappelle que connaître, c'est analyser, décomposer ce que la vie compose ; entre le geste de la vie, qui forme, et celui du savoir, qui sépare, il y a un écart que le savoir ne peut abolir. Ce que le savant rapporte n'est jamais la vie elle-même.

\begin{center}\rule{0.5\linewidth}{0.5pt}\end{center}

Faut-il alors conclure que la phrase de Sand n'éclaire rien, et qu'elle appartient à un siècle que nos œuvres auraient corrigé ? Ce serait manquer ce qu'elle a de juste et qui n'apparaît qu'à la relire jusqu'à son dernier mot : ce que Sand promet n'est pas la possession de la nature, mais celle de « son être ».

Il faut d'abord renoncer à l'image du seuil. On n'entre pas dans la nature, parce qu'on n'en est jamais sorti. C'est la leçon la plus forte de Canguilhem : l'homme de science peut construire un univers sans centre ni lieu privilégié, il n'en reste pas moins un vivant pour qui le milieu s'organise autour de lui ; la pensée qui prétend connaître la vie est une production de la vie. Il n'y a donc pas de porte à pousser, mais une séparation, instituée par le savoir lui-même, qu'il faut reconnaître. Haushofer donne à cette idée sa figure la plus saisissante : en enfermant la narratrice, le mur lui révèle qu'elle était déjà dedans. Soumise au froid, à la faim et à la fatigue comme sa vache et son chien, elle se découvre une créature parmi d'autres ; et sur l'alpage, dans le silence de la prairie, elle éprouve qu'il lui devient presque impossible de rester un moi séparé du reste du vivant. Verne, de son côté, montre que l'immersion passe par des machines : le scaphandre, les réservoirs d'air, le \emph{Nautilus} lui-même, dont la devise, \emph{Mobilis in mobili}, dit un être mobile dans l'élément mobile. La machine n'y est pas l'ennemie de la nature ; elle est, au sens de Canguilhem, le prolongement d'un vivant qui étend son milieu. Les grandes vitres du salon ne séparent pas seulement Aronnax de l'océan : elles le lui donnent à voir.

Il faut ensuite prendre au sérieux la modestie de Sand : elle n'annonce pas un secret, elle annonce « un mot », et encore « ne fût-ce qu'un mot ». Or c'est précisément ce que rapportent nos trois œuvres, qui sont elles-mêmes ce mot. Aronnax écrit un récit qu'il sait incomplet, mais il peut affirmer qu'à la question de l'Ecclésiaste deux hommes seulement ont désormais le droit de répondre, Nemo et lui. La narratrice d'Haushofer écrit sans style, sans lecteur assuré, presque sans espoir ; mais tant qu'elle écrit, elle demeure quelqu'un, et son rapport lui permet de mettre de l'ordre dans ce qu'elle a vécu, jusqu'à la mort de Lynx et du taureau, vers laquelle tout son récit s'achemine. Canguilhem, lui, ne rapporte pas de la biologie des résultats, mais une manière de penser : une philosophie qui accepte d'apprendre du vivant au lieu de lui imposer ses catégories. Le mot rapporté ne délivre pas le sens de la nature ; il maintient vivant celui qui l'écrit et le rend transmissible. C'est un usage plus modeste que l'oracle de Fontainebleau, et plus solide.

Il faut enfin donner à « la possession de son être » son sens véritable. Posséder son être, ce n'est pas s'appartenir comme on possède un bien ; Canguilhem fournit le concept qui convient. La vie, selon lui, est une activité normative : être en bonne santé, ce n'est pas se conformer à une moyenne, c'est pouvoir instituer de nouvelles normes quand le milieu change, pouvoir tomber malade et s'en relever. Se posséder, en ce sens, ce n'est pas se tenir, c'est pouvoir changer. Telle est exactement la conquête de la narratrice d'Haushofer. Les normes sociales ont disparu avec les hommes ; elle en institue d'autres, un calendrier des travaux, une répartition des réserves, une règle de conduite envers les bêtes, et ce sont ces normes qui font d'elle un sujet plutôt qu'une simple survivante. Il n'y a chez elle aucun « charme profond », aucune extase dans la forêt sacrée ; il y a le soin, qui est devenu sa manière d'être. Verne fournit le contre-exemple et, par là, la confirmation. Nemo possède la mer, l'indépendance, un empire où il ne reconnaît aucun maître ; mais il ne se possède pas, car sa haine décide de sa route jusqu'à l'hécatombe, et la nuit de la fuite Aronnax l'entend implorer Dieu de faire cesser ce qu'il a lui-même déchaîné. Aronnax, lui, ne possède rien et rapporte un récit lacunaire : c'est pourtant lui qui a changé. Le roman sépare ainsi en deux personnages ce que la phrase de Sand risquait de confondre, la conquête de la nature et la possession de soi.

\begin{center}\rule{0.5\linewidth}{0.5pt}\end{center}

La phrase de George Sand éclaire donc nos trois œuvres par son point de départ et par son terme plus que par le chemin qu'elle décrit entre les deux. Elle a raison de poser qu'il faut entrer dans la nature pour la connaître : Verne, Canguilhem et Haushofer refusent tous trois le savoir de cabinet, le vivant isolé de son milieu, le paysage regardé de loin. Elle a raison encore de tenir que ce savoir revient sur celui qui l'acquiert. Mais les mots qu'elle emploie, conquête, oracle, rapporter, sont démentis par les œuvres : la conquête y tourne à la prédation, l'oracle s'y tait parce que la nature n'a pas de sens en soi, et le retour, chez Nemo comme chez la narratrice, reste incertain ou impossible.

Ce qui subsiste n'est pas rien : c'est même l'essentiel. On ne connaît pas le vivant sans en être ; on ne rapporte de la nature qu'un mot, et ce mot, chez nos trois auteurs, est un livre ; surtout, la conquête annoncée n'était pas celle de la nature. Ce que Sand nomme la possession de son être, nos œuvres l'appellent autrement : la normativité du vivant chez Canguilhem, le soin et la règle que s'invente la narratrice d'Haushofer, la liberté intérieure qui manque à Nemo, maître de la mer et esclave de sa vengeance.

Reste la clause que Sand jette dans une incise et qui, en 1872, engageait le sort d'une forêt réelle : « pour tout homme actuellement vivant ou à naître ». Si l'expérience de la nature est due aux hommes à naître, alors elle oblige ceux qui vivent, et la conquête à entreprendre se renverse en une responsabilité à assumer.

\end{document}
